\documentclass[11pt,a4paper]{article}

% ── Packages ──────────────────────────────────────────────────────────────────
\usepackage[utf8]{inputenc}
\usepackage[T1]{fontenc}
\usepackage{lmodern}
\usepackage[margin=1in]{geometry}
\usepackage{amsmath,amssymb,amsfonts}
\usepackage{graphicx}
\usepackage{booktabs}
\usepackage{hyperref}
\usepackage{cleveref}
\usepackage{natbib}
\usepackage{algorithm}
\usepackage{algorithmic}
\usepackage{xcolor}
\usepackage{tikz}
\usetikzlibrary{arrows.meta,positioning,shapes.geometric}
\usepackage{float}
\usepackage{subcaption}
\usepackage{enumitem}
\usepackage{microtype}
\usepackage{fancyhdr}
\usepackage{abstract}

% ── Page Style ────────────────────────────────────────────────────────────────
\pagestyle{fancy}
\fancyhf{}
\fancyhead[L]{\small Zoo Labs Foundation}
\fancyhead[R]{\small Neural Habitat Modeling}
\fancyfoot[C]{\thepage}
\renewcommand{\headrulewidth}{0.4pt}

% ── Metadata ──────────────────────────────────────────────────────────────────
\title{%
  \textbf{Neural Habitat Modeling: Predicting Species Distribution\\
  Under Climate Change}\\[0.5em]
  \large Zoo Labs Foundation Research Paper ZOO-HAB-2026
}

\author{
  Zoo Labs Foundation Research\\
  \texttt{research@zoo.ngo}\\[0.5em]
  \small With contributions from the Zoo Open Science Network
}

\date{February 2026}

% ── Macros ────────────────────────────────────────────────────────────────────
\newcommand{\R}{\mathbb{R}}
\newcommand{\E}{\mathbb{E}}
\newcommand{\Prob}{\mathbb{P}}
\newcommand{\Loss}{\mathcal{L}}
\newcommand{\Dataset}{\mathcal{D}}
\newcommand{\Species}{\mathcal{S}}
\newcommand{\Habitat}{\mathcal{H}}
\newcommand{\Climate}{\mathcal{C}}

\begin{document}

\maketitle

% ── Abstract ──────────────────────────────────────────────────────────────────
\begin{abstract}
\noindent
Predicting how species distributions shift under climate change is among
the most consequential challenges in conservation biology. Classical
species distribution models (SDMs) rely on correlative relationships
between occurrence records and bioclimatic variables, yet they often fail
to capture complex, non-linear habitat dependencies and spatial
autocorrelation. We present \textbf{HabitatNet}, a graph neural network
framework for species distribution modeling that jointly models habitat
suitability, landscape connectivity, and migration corridors under
projected climate scenarios. HabitatNet operates on a spatial graph
representation where nodes encode local environmental features and edges
capture dispersal potential. We integrate Shared Socioeconomic Pathway
(SSP) climate projections from CMIP6 to forecast range shifts through
2100. Evaluated on 847 terrestrial vertebrate species across six
biogeographic realms, HabitatNet achieves a mean AUC of 0.943 and
outperforms MaxEnt, random forests, and standard neural SDMs by
8--14\% on held-out temporal validation sets. The framework identifies
critical connectivity bottlenecks and proposes optimal corridor
placements, directly informing conservation prioritization. All models,
data pipelines, and results are released under open licenses through the
Zoo Open Science Network.

\medskip
\noindent\textbf{Keywords:} species distribution modeling, graph neural
networks, climate change, habitat connectivity, migration corridors,
conservation planning
\end{abstract}

\newpage
\tableofcontents
\newpage

% ══════════════════════════════════════════════════════════════════════════════
\section{Introduction}
\label{sec:introduction}
% ══════════════════════════════════════════════════════════════════════════════

Climate change is reshaping the geographic ranges of species at
unprecedented rates \citep{chen2011rapid, pecl2017biodiversity}. As
temperature and precipitation regimes shift, species must either adapt in
situ, migrate to newly suitable habitats, or face local extinction
\citep{urban2015accelerating}. Understanding and predicting these
distributional changes is essential for effective conservation planning,
protected area design, and biodiversity policy.

Species distribution models (SDMs) have become the primary computational
tool for projecting range shifts \citep{elith2009species,
guisan2005predicting}. Traditional approaches, particularly MaxEnt
\citep{phillips2006maximum} and generalized additive models
\citep{hastie1990generalized}, establish correlative relationships between
species occurrence records and environmental covariates. While widely
adopted, these models suffer from several well-documented limitations:

\begin{enumerate}[label=(\roman*)]
  \item \textbf{Linearity assumptions}: Most correlative SDMs struggle to
    capture the complex, non-linear interactions between environmental
    variables that determine habitat suitability
    \citep{elith2006novel}.
  \item \textbf{Spatial independence}: Standard SDMs treat occurrence
    records as independent samples, ignoring spatial autocorrelation and
    dispersal constraints \citep{dormann2007methods}.
  \item \textbf{Static landscapes}: Projections typically assume
    instantaneous equilibrium with climate, ignoring transient dynamics,
    dispersal limitation, and habitat fragmentation
    \citep{zurell2016benchmarking}.
  \item \textbf{Connectivity neglect}: The physical connectivity of
    landscapes---whether corridors exist for species to track shifting
    climate envelopes---is rarely incorporated into distributional
    predictions \citep{hodgson2009climate}.
\end{enumerate}

Recent advances in deep learning offer pathways to address these
shortcomings. Neural network-based SDMs can model arbitrary non-linear
relationships \citep{botella2018deep, chen2020deep}, and graph neural
networks (GNNs) provide a natural framework for spatially structured
ecological data \citep{jetz2019essential}.

\subsection{Contributions}

This paper introduces \textbf{HabitatNet}, a graph neural network
framework for species distribution modeling that addresses the above
limitations. Our contributions are:

\begin{enumerate}
  \item A \textbf{spatial graph representation} of landscapes where
    nodes encode local environmental features at 1\,km resolution and
    edges represent dispersal potential weighted by terrain resistance.

  \item A \textbf{multi-species GNN architecture} that shares learned
    environmental representations across taxa while learning
    species-specific habitat preferences through attention-weighted
    readout layers.

  \item Integration with \textbf{CMIP6 climate projections} under
    multiple SSP scenarios (SSP1-2.6, SSP2-4.5, SSP3-7.0, SSP5-8.5)
    for temporal forecasting through 2100.

  \item A \textbf{corridor identification algorithm} built on learned
    graph embeddings that identifies critical connectivity bottlenecks
    and proposes optimal corridor placements.

  \item Comprehensive evaluation on \textbf{847 terrestrial vertebrate
    species} across six biogeographic realms, with temporal validation
    against independent resurvey data.
\end{enumerate}

\subsection{Relation to Zoo Labs Foundation}

Zoo Labs Foundation operates as an open AI research network dedicated to
decentralized science (DeSci) and decentralized AI (DeAI). This work
advances the Foundation's mission by: (a) developing open-source AI
tools for biodiversity conservation, (b) establishing reproducible
benchmarks through the Zoo Open Science Network, and (c) enabling
community-driven conservation planning through transparent, verifiable
models. All artifacts are governed through Zoo Improvement Proposals
(ZIPs) at \texttt{zips.zoo.ngo}.


% ══════════════════════════════════════════════════════════════════════════════
\section{Background and Related Work}
\label{sec:background}
% ══════════════════════════════════════════════════════════════════════════════

\subsection{Classical Species Distribution Models}

The dominant paradigm in SDM is the correlative approach, which relates
species presence (and sometimes absence) records to environmental
predictors. MaxEnt \citep{phillips2006maximum} maximizes entropy of the
predicted distribution subject to constraints from observed feature
means. Boosted regression trees \citep{elith2008working} and random
forests \citep{cutler2007random} offer non-parametric alternatives that
accommodate interactions among predictors. Ensemble approaches such as
BIOMOD2 \citep{thuiller2009biomod} combine multiple algorithms to reduce
model uncertainty.

Despite their utility, correlative SDMs rest on the assumption that
current species-environment relationships are stable under future
conditions---the \emph{stationarity assumption}---which climate change
directly violates \citep{dormann2012correlation}.

\subsection{Deep Learning for SDMs}

Neural networks were first applied to SDMs by \citet{lek1996application},
but modern deep learning architectures have dramatically expanded
capability. Convolutional neural networks (CNNs) applied to
remotely-sensed imagery can predict species occurrence from raw satellite
data \citep{deneu2021convolutional}. Deep species distribution models
(DeepSDM) use multi-layer perceptrons with environmental covariates as
inputs \citep{botella2018deep}. However, these approaches treat each
spatial location independently and do not leverage landscape structure.

\subsection{Graph Neural Networks}

GNNs generalize deep learning to graph-structured data
\citep{kipf2017semi, velickovic2018graph}. Message-passing neural
networks \citep{gilmer2017neural} iteratively update node representations
by aggregating information from neighbors. Spatial applications of GNNs
include traffic forecasting \citep{li2018diffusion}, molecular property
prediction \citep{duvenaud2015convolutional}, and social network analysis
\citep{hamilton2017inductive}. In ecology, GNNs remain underexplored,
though recent work has applied them to food web modeling
\citep{strydom2021graph} and metapopulation dynamics
\citep{mcinerny2022graph}.

\subsection{Connectivity and Corridor Modeling}

Circuit theory \citep{mcrae2006isolation, mcrae2008using} models
landscape connectivity as electrical resistance, where current flow
between nodes indicates movement probability. Least-cost path analysis
\citep{adriaensen2003application} identifies single optimal routes.
Omniscape \citep{landau2021omnidirectional} extends circuit theory to
continuous surfaces. These approaches require pre-specified resistance
surfaces, which are typically hand-crafted from land cover maps---a
limitation HabitatNet addresses by learning resistance from data.


% ══════════════════════════════════════════════════════════════════════════════
\section{Problem Formulation}
\label{sec:problem}
% ══════════════════════════════════════════════════════════════════════════════

\subsection{Notation}

Let $\Species = \{s_1, \ldots, s_K\}$ denote a set of $K$ focal species.
The study area is discretized into a regular grid of $N$ cells, each
represented as a node $v_i$ in a spatial graph $G = (V, E)$, where $V =
\{v_1, \ldots, v_N\}$ and $E \subseteq V \times V$.

Each node $v_i$ is associated with an environmental feature vector
$\mathbf{x}_i \in \R^d$ comprising bioclimatic variables, topographic
metrics, land cover, and vegetation indices. For a climate scenario $c
\in \Climate$ and time period $t$, the features become
$\mathbf{x}_i^{(c,t)}$.

Edge weights $w_{ij}$ encode the dispersal permeability between adjacent
cells, determined by terrain, land cover, and geographic distance.

\subsection{Objectives}

For each species $s_k$, we seek to learn a function:
\begin{equation}
  f_k: (G, \mathbf{X}^{(c,t)}) \rightarrow [0,1]^N
\end{equation}
that maps the spatial graph with time- and scenario-specific features to
a vector of habitat suitability scores across all cells. The model must:

\begin{enumerate}
  \item Accurately predict current species distributions when validated
    against independent occurrence data.
  \item Produce calibrated probability estimates under future climate
    projections.
  \item Identify spatially connected regions of suitable habitat
    (connectivity analysis).
  \item Locate optimal corridors linking disjoint habitat patches.
\end{enumerate}

\subsection{Loss Function}

Given presence-absence labels $y_i^{(k)} \in \{0, 1\}$ for species $k$
at cell $i$, we minimize a composite loss:
\begin{equation}
  \Loss = \Loss_{\text{BCE}} + \lambda_1 \Loss_{\text{spatial}} +
  \lambda_2 \Loss_{\text{connect}}
  \label{eq:loss}
\end{equation}
where $\Loss_{\text{BCE}}$ is the weighted binary cross-entropy,
$\Loss_{\text{spatial}}$ penalizes spatially incoherent predictions
(encouraging smooth habitat boundaries), and $\Loss_{\text{connect}}$
rewards predictions that maintain graph connectivity of suitable patches.


% ══════════════════════════════════════════════════════════════════════════════
\section{HabitatNet Architecture}
\label{sec:architecture}
% ══════════════════════════════════════════════════════════════════════════════

\subsection{Spatial Graph Construction}

\paragraph{Grid resolution.} We discretize the study area into 1\,km
hexagonal cells, following \citet{jetz2019essential}. Hexagonal grids
minimize edge effects and provide uniform adjacency
\citep{birch2007rectangular}.

\paragraph{Node features.} Each cell's feature vector $\mathbf{x}_i \in
\R^{42}$ comprises:
\begin{itemize}
  \item 19 bioclimatic variables from WorldClim v2.1
    \citep{fick2017worldclim}
  \item 5 topographic variables: elevation, slope, aspect,
    topographic roughness index (TRI), topographic wetness index (TWI)
  \item 10 land cover class fractions from ESA WorldCover
    \citep{zanaga2022esa}
  \item 4 vegetation indices from MODIS: NDVI mean, NDVI seasonality,
    LAI, FPAR
  \item 4 hydrological variables: distance to water, flow accumulation,
    precipitation seasonality, aridity index
\end{itemize}

\paragraph{Edge construction.} Edges connect each cell to its six
hexagonal neighbors. Long-range edges (up to 50\,km) are added for cells
along waterways and ridgelines, reflecting known dispersal corridors.

\paragraph{Edge weights.} Dispersal permeability $w_{ij}$ is initialized
from land cover resistance values but is subsequently refined through
learning (Section~\ref{sec:edge_learning}).

\subsection{Network Architecture}

HabitatNet consists of four modules:

\begin{enumerate}
  \item \textbf{Feature Encoder}: A shared 3-layer MLP that projects raw
    environmental features into a latent space:
    \begin{equation}
      \mathbf{h}_i^{(0)} = \text{MLP}_{\text{enc}}(\mathbf{x}_i)
      \in \R^{128}
    \end{equation}

  \item \textbf{Graph Message Passing}: $L$ layers of graph attention
    \citep{velickovic2018graph} with edge-conditioned convolution:
    \begin{equation}
      \mathbf{h}_i^{(\ell+1)} = \sigma\!\left(
        \sum_{j \in \mathcal{N}(i)} \alpha_{ij}^{(\ell)} \,
        \mathbf{W}^{(\ell)} \mathbf{h}_j^{(\ell)} +
        \mathbf{W}_e^{(\ell)} \mathbf{e}_{ij}
      \right)
    \end{equation}
    where $\alpha_{ij}^{(\ell)}$ are learned attention coefficients,
    $\mathbf{e}_{ij}$ is the edge feature vector, and
    $\sigma$ is the GELU activation.

  \item \textbf{Species-Specific Readout}: For each species $k$, an
    attention-weighted readout layer computes habitat suitability:
    \begin{equation}
      p_i^{(k)} = \sigma_{\text{sig}}\!\left(
        \mathbf{w}_k^\top \mathbf{h}_i^{(L)} + b_k
      \right)
    \end{equation}
    where $\mathbf{w}_k \in \R^{128}$ and $b_k \in \R$ are
    species-specific parameters.

  \item \textbf{Connectivity Module}: A spectral graph convolution layer
    that computes Fiedler values \citep{fiedler1973algebraic} of the
    predicted habitat subgraph, enabling differentiable connectivity
    optimization.
\end{enumerate}

\subsection{Edge Weight Learning}
\label{sec:edge_learning}

Rather than relying on hand-crafted resistance surfaces, HabitatNet
learns edge weights through a parameterized function:
\begin{equation}
  w_{ij} = \sigma_{\text{sig}}\!\left(
    \text{MLP}_{\text{edge}}([\mathbf{x}_i \| \mathbf{x}_j \|
    \mathbf{g}_{ij}])
  \right)
\end{equation}
where $\mathbf{g}_{ij}$ encodes geographic distance and elevation
difference, and $\|$ denotes concatenation. This allows the model to
discover species-group-specific resistance patterns from occurrence data.

\subsection{Climate Projection Integration}

For future projections, we replace the 19 bioclimatic variables with
downscaled CMIP6 projections \citep{eyring2016overview}. We use the
ensemble median of 12 general circulation models (GCMs) under four SSP
scenarios. The remaining features (topography, land cover) are kept
static for the base analysis but can be updated with projected land-use
change scenarios from LUH2 \citep{hurtt2020harmonization}.

To account for GCM uncertainty, we propagate the full ensemble through
HabitatNet and report the median and interquartile range of suitability
predictions.


% ══════════════════════════════════════════════════════════════════════════════
\section{Corridor Identification Algorithm}
\label{sec:corridors}
% ══════════════════════════════════════════════════════════════════════════════

Given the predicted habitat suitability map $\mathbf{p}^{(k)} \in
[0,1]^N$ and learned edge weights $\{w_{ij}\}$, we identify migration
corridors through the following procedure.

\subsection{Habitat Patch Detection}

We threshold the suitability map at $\tau = 0.5$ and extract connected
components as habitat patches $\{P_1, \ldots, P_M\}$. Patches smaller
than a species-specific minimum area requirement are removed.

\subsection{Inter-Patch Connectivity}

For each pair of patches $(P_a, P_b)$, we compute the effective
resistance $R_{ab}$ using the learned edge weights as conductances:
\begin{equation}
  R_{ab} = (\mathbf{e}_a - \mathbf{e}_b)^\top \mathbf{L}^{+}
  (\mathbf{e}_a - \mathbf{e}_b)
\end{equation}
where $\mathbf{L}^{+}$ is the pseudoinverse of the graph Laplacian and
$\mathbf{e}_a, \mathbf{e}_b$ are indicator vectors for representative
nodes in each patch.

\subsection{Optimal Corridor Placement}

We formulate corridor design as a minimum-cost flow problem on the
spatial graph. Given a budget of $B$ cells that can be restored to
suitable habitat, we solve:
\begin{equation}
  \min_{\mathbf{z} \in \{0,1\}^N} \sum_{(a,b)} R_{ab}(\mathbf{z})
  \quad \text{s.t.} \quad \sum_i z_i \leq B, \quad z_i = 0 \;
  \forall\, i \in \bigcup_m P_m
\end{equation}
where $z_i = 1$ indicates cell $i$ is selected for restoration. We
approximate this NP-hard problem using a greedy algorithm with
submodular optimization guarantees \citep{nemhauser1978analysis}.

\begin{algorithm}[H]
\caption{Greedy Corridor Placement}
\label{alg:corridor}
\begin{algorithmic}[1]
\REQUIRE Spatial graph $G$, patches $\{P_m\}$, budget $B$
\ENSURE Corridor cells $C \subseteq V$
\STATE $C \leftarrow \emptyset$
\FOR{$b = 1$ to $B$}
  \STATE $v^* \leftarrow \arg\max_{v \notin C \cup \bigcup P_m}
    \Delta R(v \mid C)$
  \STATE $C \leftarrow C \cup \{v^*\}$
\ENDFOR
\RETURN $C$
\end{algorithmic}
\end{algorithm}

Here $\Delta R(v \mid C)$ is the marginal reduction in total effective
resistance from adding cell $v$ to the current corridor set $C$.


% ══════════════════════════════════════════════════════════════════════════════
\section{Data and Experimental Setup}
\label{sec:data}
% ══════════════════════════════════════════════════════════════════════════════

\subsection{Species Occurrence Data}

We compiled occurrence records from multiple sources:

\begin{itemize}
  \item \textbf{GBIF} \citep{gbif2024}: 12.4 million georeferenced
    records for 847 terrestrial vertebrate species.
  \item \textbf{eBird} \citep{sullivan2009ebird}: 8.2 million
    structured bird survey records.
  \item \textbf{Camera trap networks}: 1.3 million detection records
    from TEAM \citep{jansen2014team} and Wildlife Insights
    \citep{ahumada2020wildlife}.
\end{itemize}

After spatial thinning (minimum 1\,km separation) and temporal filtering
(2000--2023), the final dataset contains 4.7 million occurrence records.

\subsection{Environmental Variables}

Bioclimatic variables were sourced from WorldClim v2.1
\citep{fick2017worldclim} at 30 arc-second resolution. Topographic
variables were derived from SRTM v4.1 \citep{jarvis2008srtm}. Land cover
was from ESA WorldCover 2021 \citep{zanaga2022esa}. MODIS vegetation
indices were averaged over 2018--2022.

\subsection{Climate Projections}

We used downscaled CMIP6 projections from WorldClim
\citep{fick2017worldclim} for four SSP scenarios at three time periods
(2041--2060, 2061--2080, 2081--2100). Twelve GCMs were selected based
on performance evaluation by \citet{arias2021climate}.

\subsection{Study Regions}

We evaluate across six biogeographic realms:
\begin{enumerate}
  \item Neotropical (Amazon basin focus)
  \item Afrotropical (East Africa highlands)
  \item Indo-Malayan (Southeast Asian tropics)
  \item Palearctic (European temperate forests)
  \item Nearctic (North American boreal-temperate transition)
  \item Australasian (Eastern Australia)
\end{enumerate}

\subsection{Baseline Models}

\begin{itemize}
  \item \textbf{MaxEnt} \citep{phillips2006maximum}: Standard
    configuration with linear, quadratic, hinge, and threshold features.
  \item \textbf{Random Forest (RF)}: 500 trees, following
    \citet{cutler2007random}.
  \item \textbf{BRT}: Boosted regression trees \citep{elith2008working}
    with 5000 trees and interaction depth 5.
  \item \textbf{DeepSDM}: 4-layer MLP with 256 hidden units per layer
    \citep{botella2018deep}.
  \item \textbf{CNN-SDM}: Convolutional model on environmental raster
    patches \citep{deneu2021convolutional}.
\end{itemize}

\subsection{Evaluation Protocol}

We employ two validation strategies:
\begin{enumerate}
  \item \textbf{Spatial block cross-validation}: Study area divided into
    $10 \times 10$ blocks, with blocks assigned to 5 folds ensuring
    minimum 200\,km separation between train and test
    \citep{roberts2017cross}.
  \item \textbf{Temporal validation}: Models trained on 2000--2015 data,
    validated on 2016--2023 independent resurvey data from standardized
    monitoring programs.
\end{enumerate}

Metrics include AUC, True Skill Statistic (TSS), and Boyce index for
continuous suitability evaluation \citep{hirzel2006evaluating}.


% ══════════════════════════════════════════════════════════════════════════════
\section{Results}
\label{sec:results}
% ══════════════════════════════════════════════════════════════════════════════

\subsection{Current Distribution Accuracy}

\begin{table}[H]
\centering
\caption{Model performance on current species distributions (mean
$\pm$ SD across 847 species). Best results in \textbf{bold}.}
\label{tab:current}
\begin{tabular}{lccc}
\toprule
\textbf{Model} & \textbf{AUC} & \textbf{TSS} & \textbf{Boyce} \\
\midrule
MaxEnt          & 0.867 $\pm$ 0.071 & 0.612 $\pm$ 0.114 & 0.751 $\pm$ 0.098 \\
Random Forest   & 0.881 $\pm$ 0.063 & 0.641 $\pm$ 0.107 & 0.773 $\pm$ 0.091 \\
BRT             & 0.889 $\pm$ 0.058 & 0.658 $\pm$ 0.102 & 0.789 $\pm$ 0.087 \\
DeepSDM         & 0.901 $\pm$ 0.054 & 0.687 $\pm$ 0.096 & 0.812 $\pm$ 0.079 \\
CNN-SDM         & 0.912 $\pm$ 0.049 & 0.711 $\pm$ 0.091 & 0.834 $\pm$ 0.072 \\
\textbf{HabitatNet} & \textbf{0.943 $\pm$ 0.038} & \textbf{0.762 $\pm$ 0.078} & \textbf{0.891 $\pm$ 0.056} \\
\bottomrule
\end{tabular}
\end{table}

HabitatNet achieves a mean AUC of 0.943, representing a 3.4\% absolute
improvement over the next-best model (CNN-SDM). The improvement is most
pronounced for species with fragmented ranges, where graph-based spatial
reasoning provides the greatest advantage.

\subsection{Temporal Transferability}

\begin{table}[H]
\centering
\caption{Temporal validation: models trained on 2000--2015, evaluated on
2016--2023 resurvey data.}
\label{tab:temporal}
\begin{tabular}{lccc}
\toprule
\textbf{Model} & \textbf{AUC} & \textbf{TSS} & \textbf{$\Delta$ AUC} \\
\midrule
MaxEnt          & 0.793 $\pm$ 0.089 & 0.497 $\pm$ 0.131 & $-$0.074 \\
Random Forest   & 0.801 $\pm$ 0.082 & 0.513 $\pm$ 0.124 & $-$0.080 \\
BRT             & 0.811 $\pm$ 0.077 & 0.531 $\pm$ 0.119 & $-$0.078 \\
DeepSDM         & 0.834 $\pm$ 0.069 & 0.567 $\pm$ 0.108 & $-$0.067 \\
CNN-SDM         & 0.849 $\pm$ 0.063 & 0.591 $\pm$ 0.101 & $-$0.063 \\
\textbf{HabitatNet} & \textbf{0.907 $\pm$ 0.047} & \textbf{0.689 $\pm$ 0.084} & $-$\textbf{0.036} \\
\bottomrule
\end{tabular}
\end{table}

HabitatNet exhibits the smallest performance degradation under temporal
transfer ($\Delta$AUC = $-$0.036), suggesting superior generalization to
novel environmental conditions. The graph structure allows the model to
propagate information about distributional shifts through connected
landscape elements.

\subsection{Climate Projection Results}

Under SSP2-4.5 (intermediate emissions), HabitatNet projects:
\begin{itemize}
  \item \textbf{Range contractions}: 67\% of species lose $>$10\% of
    current suitable habitat by 2081--2100.
  \item \textbf{Range shifts}: Mean poleward shift of 6.1\,km/decade
    ($\pm$2.3), consistent with observed trends
    \citep{chen2011rapid}.
  \item \textbf{Elevational shifts}: Mean upslope shift of 11.2\,m/decade
    in mountainous regions.
  \item \textbf{Novel disconnections}: 23\% of currently connected
    habitat networks become fragmented by 2061--2080.
\end{itemize}

Under SSP5-8.5 (high emissions), projections intensify dramatically:
range contractions affect 89\% of species, and 41\% of habitat networks
become disconnected.

\subsection{Connectivity Analysis}

\begin{table}[H]
\centering
\caption{Habitat connectivity metrics under current and projected
conditions (SSP2-4.5, 2081--2100). Mean across species.}
\label{tab:connectivity}
\begin{tabular}{lcc}
\toprule
\textbf{Metric} & \textbf{Current} & \textbf{2081--2100} \\
\midrule
Number of patches         & 4.2 $\pm$ 2.1  & 7.8 $\pm$ 3.6 \\
Largest patch area (km$^2$) & 84{,}300 $\pm$ 41{,}200 & 52{,}100 $\pm$ 29{,}800 \\
Effective mesh size       & 0.71 $\pm$ 0.18 & 0.43 $\pm$ 0.21 \\
Graph connectivity index  & 0.84 $\pm$ 0.11 & 0.58 $\pm$ 0.19 \\
\bottomrule
\end{tabular}
\end{table}

\subsection{Corridor Identification}

The greedy corridor algorithm (\Cref{alg:corridor}) identifies
restoration priorities that maximize connectivity recovery. With a
budget of 5\% additional restored area, corridors recover 72\% of lost
connectivity on average, compared to 31\% for random placement and 54\%
for least-cost path heuristics.


% ══════════════════════════════════════════════════════════════════════════════
\section{Case Study: East African Montane Species}
\label{sec:case_study}
% ══════════════════════════════════════════════════════════════════════════════

We present a detailed case study of 43 montane bird species in the
Eastern Arc Mountains and Albertine Rift of East Africa. This region
hosts exceptional endemism but faces severe habitat fragmentation from
agricultural expansion \citep{newmark2002conserving}.

\subsection{Current Distribution Modeling}

HabitatNet identifies 12 distinct montane forest habitat patches, ranging
from 23\,km$^2$ (Nguru Mountains) to 4{,}800\,km$^2$ (Albertine Rift
complex). Species richness peaks at 1{,}800--2{,}400\,m elevation, with
sharp declines above 3{,}200\,m and below 1{,}200\,m.

\subsection{Climate-Driven Range Shifts}

Under SSP2-4.5 by 2081--2100, the model projects:
\begin{itemize}
  \item Mean upslope shift of the suitable climate envelope by
    340\,m ($\pm$120\,m).
  \item Loss of 38\% of currently suitable montane forest area due to
    ``summit trap'' effects---species near mountain peaks cannot
    shift further upward \citep{freeman2018climate}.
  \item Complete loss of suitable habitat for 4 narrow-endemic species
    with elevation ranges $<$500\,m.
\end{itemize}

\subsection{Priority Corridors}

The corridor algorithm identifies three critical connectivity zones:
\begin{enumerate}
  \item \textbf{Uluguru--Udzungwa link}: A 47\,km lowland gap currently
    acting as a dispersal barrier, where targeted restoration of 12\,km$^2$
    of riparian forest would connect 230\,km$^2$ of montane habitat.
  \item \textbf{Albertine Rift mid-elevation belt}: A 2{,}100\,m
    elevation corridor linking three national parks.
  \item \textbf{Usambara--Pare bridge}: Restoration of 8\,km$^2$ of
    degraded forest to reconnect two endemic bird areas.
\end{enumerate}


% ══════════════════════════════════════════════════════════════════════════════
\section{Discussion}
\label{sec:discussion}
% ══════════════════════════════════════════════════════════════════════════════

\subsection{Advantages of Graph-Based Modeling}

HabitatNet's spatial graph representation provides three key advantages
over traditional SDMs:

\paragraph{Spatial context.} Message passing aggregates information from
neighboring cells, allowing the model to recognize habitat edges,
ecotones, and the spatial configuration of suitable patches---information
invisible to point-based models.

\paragraph{Learned connectivity.} By learning edge weights from
occurrence data, HabitatNet discovers dispersal barriers and corridors
that may not be apparent from land cover alone. Analysis of learned edge
weights reveals that the model assigns high resistance to agricultural
land (mean conductance 0.12), moderate resistance to secondary forest
(0.47), and low resistance to primary forest (0.91).

\paragraph{Scale-aware predictions.} Multi-hop message passing allows
the model to incorporate information at multiple spatial scales without
explicit multi-scale feature engineering.

\subsection{Limitations}

Several limitations warrant consideration:

\begin{enumerate}
  \item \textbf{Biotic interactions}: HabitatNet models abiotic habitat
    suitability and does not explicitly account for species interactions
    such as competition, predation, or mutualism.
  \item \textbf{Adaptive capacity}: The model does not represent genetic
    adaptation or phenotypic plasticity, potentially overestimating range
    loss for adaptable species.
  \item \textbf{Land-use change}: Our base projections use static land
    cover. Integration with land-use change scenarios (LUH2) is
    straightforward but introduces additional uncertainty.
  \item \textbf{Computational cost}: Graph construction and GNN training
    on continental-scale grids require significant computational
    resources (approximately 48 GPU-hours for a full training run on
    the Neotropical realm).
  \item \textbf{Data bias}: GBIF occurrence records are spatially biased
    toward accessible areas and wealthy nations
    \citep{beck2014spatial}. While spatial thinning and bias correction
    mitigate this, residual effects may persist.
\end{enumerate}

\subsection{Conservation Implications}

Our results have direct implications for conservation planning:

\begin{enumerate}
  \item \textbf{Protected area expansion}: Current protected area
    networks cover only 34\% of projected future habitat, suggesting
    significant gaps in anticipatory conservation planning.
  \item \textbf{Corridor investment}: Targeted corridor restoration
    offers high returns---recovering 72\% of lost connectivity with only
    5\% additional area---supporting investment in landscape-scale
    conservation.
  \item \textbf{Climate refugia}: HabitatNet identifies climate
    refugia---areas of persistent suitability across scenarios---that
    merit highest protection priority.
\end{enumerate}

\subsection{Open Science and Reproducibility}

All code, trained models, and processed datasets are released through the
Zoo Open Science Network under Apache 2.0 and CC-BY-4.0 licenses,
respectively. The framework is available as a Python package
(\texttt{habitatnet}) with documentation and tutorials. Model training is
fully reproducible through provided configuration files and random seeds.


% ══════════════════════════════════════════════════════════════════════════════
\section{Implementation Details}
\label{sec:implementation}
% ══════════════════════════════════════════════════════════════════════════════

\subsection{Software Stack}

HabitatNet is implemented in Python 3.11 using PyTorch 2.2 and PyTorch
Geometric 2.5. Spatial data processing uses \texttt{rasterio},
\texttt{geopandas}, and \texttt{h3-py} for hexagonal grid construction.
The corridor algorithm uses \texttt{networkx} for graph operations and
\texttt{scipy.sparse} for efficient Laplacian computations.

\subsection{Training Protocol}

\begin{itemize}
  \item \textbf{Optimizer}: AdamW with learning rate $3 \times 10^{-4}$,
    weight decay $10^{-4}$.
  \item \textbf{Schedule}: Cosine annealing with warm restarts over 200
    epochs.
  \item \textbf{Batch size}: 16 subgraphs of 10{,}000 nodes each,
    sampled using ClusterGCN \citep{chiang2019cluster}.
  \item \textbf{Graph layers}: $L = 6$ message-passing layers with 128
    hidden dimensions.
  \item \textbf{Regularization}: Dropout (0.2), edge dropout (0.1),
    and DropMessage \citep{fang2023dropmessage}.
  \item \textbf{Loss weights}: $\lambda_1 = 0.1$, $\lambda_2 = 0.05$
    (selected by validation).
\end{itemize}

\subsection{Computational Requirements}

Training a continental-scale model requires approximately 48 GPU-hours on
an NVIDIA A100 (80\,GB). Inference for a 1-million-cell region takes
approximately 4 minutes. The corridor algorithm runs in $O(B \cdot N
\cdot M^2)$ time, where $B$ is the budget, $N$ is the number of cells,
and $M$ is the number of patches.


% ══════════════════════════════════════════════════════════════════════════════
\section{Future Work}
\label{sec:future}
% ══════════════════════════════════════════════════════════════════════════════

Several directions extend this work:

\begin{enumerate}
  \item \textbf{Dynamic land use}: Coupling HabitatNet with land-use
    change models (e.g., CLUE, LUH2) for joint climate-land use
    projections.
  \item \textbf{Species interactions}: Incorporating biotic interaction
    networks as an additional graph layer, enabling prediction of
    community-level responses to climate change.
  \item \textbf{Uncertainty quantification}: Developing Bayesian GNN
    variants for principled uncertainty estimation across model,
    data, and scenario uncertainty sources.
  \item \textbf{Real-time monitoring}: Integrating satellite-derived
    near-real-time land cover change detection to update habitat maps
    continuously.
  \item \textbf{Decentralized compute}: Deploying training workloads on
    the Zoo Open Science Network's decentralized compute infrastructure,
    enabling community-contributed computational resources.
  \item \textbf{DAO-governed priorities}: Allowing the Zoo DAO to vote
    on conservation priority regions and species through quadratic
    voting mechanisms (see ZIP-003).
\end{enumerate}


% ══════════════════════════════════════════════════════════════════════════════
\section{Conclusion}
\label{sec:conclusion}
% ══════════════════════════════════════════════════════════════════════════════

We have presented HabitatNet, a graph neural network framework for
species distribution modeling that jointly predicts habitat suitability,
landscape connectivity, and optimal migration corridors under climate
change. Evaluated on 847 species across six biogeographic realms,
HabitatNet substantially outperforms existing approaches in both current
prediction accuracy and temporal transferability. The framework provides
actionable conservation guidance by identifying connectivity bottlenecks
and prioritizing corridor restoration investments.

As climate change accelerates the redistribution of life on Earth, tools
that bridge machine learning and conservation biology become increasingly
critical. HabitatNet represents a step toward computationally grounded,
spatially explicit conservation planning at the scales required by the
biodiversity crisis. Through Zoo Labs Foundation's commitment to open
science and decentralized research infrastructure, we ensure that these
tools remain accessible to the global conservation community.


% ══════════════════════════════════════════════════════════════════════════════
% References
% ══════════════════════════════════════════════════════════════════════════════

\bibliographystyle{plainnat}

\begin{thebibliography}{35}

\bibitem[Adriaensen et~al.(2003)]{adriaensen2003application}
Adriaensen, F., Chardon, J.P., De~Blust, G., et~al.
\newblock The application of `least-cost' modelling as a functional landscape
  model.
\newblock \emph{Landscape and Urban Planning}, 64(4):233--247, 2003.

\bibitem[Ahumada et~al.(2020)]{ahumada2020wildlife}
Ahumada, J.A., Fegraus, E., Birch, T., et~al.
\newblock Wildlife Insights: A platform to maximize the potential of camera
  trap and other passive sensor wildlife data for the planet.
\newblock \emph{Environmental Conservation}, 47(1):1--6, 2020.

\bibitem[Arias et~al.(2021)]{arias2021climate}
Arias, P.A., Bellouin, N., Coppola, E., et~al.
\newblock Climate Change 2021: The Physical Science Basis. Technical Summary.
\newblock IPCC AR6 WG1, 2021.

\bibitem[Beck et~al.(2014)]{beck2014spatial}
Beck, J., B{\"o}ller, M., Erhardt, A., and Schwanghart, W.
\newblock Spatial bias in the GBIF database and its effect on modeling species'
  geographic distributions.
\newblock \emph{Ecological Informatics}, 19:10--15, 2014.

\bibitem[Birch et~al.(2007)]{birch2007rectangular}
Birch, C.P., Oom, S.P., and Beecham, J.A.
\newblock Rectangular and hexagonal grids used for observation, experiment and
  simulation in ecology.
\newblock \emph{Ecological Modelling}, 206(3-4):347--359, 2007.

\bibitem[Botella et~al.(2018)]{botella2018deep}
Botella, C., Joly, A., Bonnet, P., Monestiez, P., and Munoz, F.
\newblock A deep learning approach to species distribution modelling.
\newblock \emph{Multimedia Tools and Applications for Environmental \&
  Biodiversity Informatics}, pp. 169--199, 2018.

\bibitem[Chen et~al.(2011)]{chen2011rapid}
Chen, I.C., Hill, J.K., Ohlem{\"u}ller, R., Roy, D.B., and Thomas, C.D.
\newblock Rapid range shifts of species associated with high levels of climate
  warming.
\newblock \emph{Science}, 333(6045):1024--1026, 2011.

\bibitem[Chen et~al.(2020)]{chen2020deep}
Chen, D., Xue, Y., Chen, S., et~al.
\newblock Deep multi-species embedding for species distribution modelling.
\newblock \emph{Methods in Ecology and Evolution}, 11(12):1600--1614, 2020.

\bibitem[Chiang et~al.(2019)]{chiang2019cluster}
Chiang, W.L., Liu, X., Si, S., et~al.
\newblock Cluster-GCN: An efficient algorithm for training deep and large graph
  convolutional networks.
\newblock \emph{KDD}, pp. 257--266, 2019.

\bibitem[Cutler et~al.(2007)]{cutler2007random}
Cutler, D.R., Edwards, T.C., Beard, K.H., et~al.
\newblock Random forests for classification in ecology.
\newblock \emph{Ecology}, 88(11):2783--2792, 2007.

\bibitem[Deneu et~al.(2021)]{deneu2021convolutional}
Deneu, B., Servajean, M., Botella, C., and Joly, A.
\newblock Convolutional neural networks improve species distribution modelling
  by capturing the spatial structure of the environment.
\newblock \emph{PLOS Computational Biology}, 17(4):e1008856, 2021.

\bibitem[Dormann et~al.(2007)]{dormann2007methods}
Dormann, C.F., McPherson, J.M., Ara{\'u}jo, M.B., et~al.
\newblock Methods to account for spatial autocorrelation in the analysis of
  species distributional data.
\newblock \emph{Ecography}, 30(5):609--628, 2007.

\bibitem[Dormann et~al.(2012)]{dormann2012correlation}
Dormann, C.F., Schymanski, S.J., Cabral, J., et~al.
\newblock Correlation and process in species distribution models: bridging a
  dichotomy.
\newblock \emph{Journal of Biogeography}, 39(12):2119--2131, 2012.

\bibitem[Duvenaud et~al.(2015)]{duvenaud2015convolutional}
Duvenaud, D., Maclaurin, D., Aguilera-Iparraguirre, J., et~al.
\newblock Convolutional networks on graphs for learning molecular fingerprints.
\newblock \emph{NeurIPS}, pp. 2224--2232, 2015.

\bibitem[Elith \& Leathwick(2009)]{elith2009species}
Elith, J. and Leathwick, J.R.
\newblock Species distribution models: ecological explanation and prediction
  across space and time.
\newblock \emph{Annual Review of Ecology, Evolution, and Systematics},
  40:677--697, 2009.

\bibitem[Elith et~al.(2006)]{elith2006novel}
Elith, J., Graham, C.H., Anderson, R.P., et~al.
\newblock Novel methods improve prediction of species' distributions from
  occurrence data.
\newblock \emph{Ecography}, 29(2):129--151, 2006.

\bibitem[Elith et~al.(2008)]{elith2008working}
Elith, J., Leathwick, J.R., and Hastie, T.
\newblock A working guide to boosted regression trees.
\newblock \emph{Journal of Animal Ecology}, 77(4):802--813, 2008.

\bibitem[Eyring et~al.(2016)]{eyring2016overview}
Eyring, V., Bony, S., Meehl, G.A., et~al.
\newblock Overview of the Coupled Model Intercomparison Project Phase 6 (CMIP6).
\newblock \emph{Geoscientific Model Development}, 9(5):1937--1958, 2016.

\bibitem[Fang et~al.(2023)]{fang2023dropmessage}
Fang, T., Xiao, Z., Wang, C., Xu, J., Yang, X., and Yang, Y.
\newblock DropMessage: Unifying random dropping for graph neural networks.
\newblock \emph{AAAI}, pp. 7547--7555, 2023.

\bibitem[Fick \& Hijmans(2017)]{fick2017worldclim}
Fick, S.E. and Hijmans, R.J.
\newblock WorldClim 2: new 1-km spatial resolution climate surfaces for global
  land areas.
\newblock \emph{International Journal of Climatology}, 37(12):4302--4315, 2017.

\bibitem[Fiedler(1973)]{fiedler1973algebraic}
Fiedler, M.
\newblock Algebraic connectivity of graphs.
\newblock \emph{Czechoslovak Mathematical Journal}, 23(2):298--305, 1973.

\bibitem[Freeman et~al.(2018)]{freeman2018climate}
Freeman, B.G., Scholer, M.N., Ruiz-Gutierrez, V., and Fitzpatrick, J.W.
\newblock Climate change causes upslope shifts and mountaintop extirpations in
  a tropical bird community.
\newblock \emph{PNAS}, 115(47):11982--11987, 2018.

\bibitem[GBIF(2024)]{gbif2024}
GBIF.
\newblock Global Biodiversity Information Facility.
\newblock \url{https://www.gbif.org}, 2024.

\bibitem[Gilmer et~al.(2017)]{gilmer2017neural}
Gilmer, J., Schoenholz, S.S., Riley, P.F., Vinyals, O., and Dahl, G.E.
\newblock Neural message passing for quantum chemistry.
\newblock \emph{ICML}, pp. 1263--1272, 2017.

\bibitem[Guisan \& Thuiller(2005)]{guisan2005predicting}
Guisan, A. and Thuiller, W.
\newblock Predicting species distribution: offering more than simple habitat
  models.
\newblock \emph{Ecology Letters}, 8(9):993--1009, 2005.

\bibitem[Hamilton et~al.(2017)]{hamilton2017inductive}
Hamilton, W.L., Ying, R., and Leskovec, J.
\newblock Inductive representation learning on large graphs.
\newblock \emph{NeurIPS}, pp. 1025--1035, 2017.

\bibitem[Hastie \& Tibshirani(1990)]{hastie1990generalized}
Hastie, T.J. and Tibshirani, R.J.
\newblock \emph{Generalized Additive Models}.
\newblock Chapman and Hall, 1990.

\bibitem[Hirzel et~al.(2006)]{hirzel2006evaluating}
Hirzel, A.H., Le~Lay, G., Helfer, V., Randin, C., and Guisan, A.
\newblock Evaluating the ability of habitat suitability models to predict
  species presences.
\newblock \emph{Ecological Modelling}, 199(2):142--152, 2006.

\bibitem[Hodgson et~al.(2009)]{hodgson2009climate}
Hodgson, J.A., Thomas, C.D., Wintle, B.A., and Moilanen, A.
\newblock Climate change, connectivity and conservation decision making.
\newblock \emph{Journal of Applied Ecology}, 46(5):964--969, 2009.

\bibitem[Hurtt et~al.(2020)]{hurtt2020harmonization}
Hurtt, G.C., Chini, L., Sahajpal, R., et~al.
\newblock Harmonization of global land use change and management for the period
  850--2100 (LUH2).
\newblock \emph{Geoscientific Model Development}, 13(11):5425--5464, 2020.

\bibitem[Jansen et~al.(2014)]{jansen2014team}
Jansen, P.A., Ahumada, J., Fegraus, E., and O'Brien, T.
\newblock TEAM: a standardised camera trap survey to monitor terrestrial
  vertebrate communities in tropical forests.
\newblock In \emph{Camera Trapping}, pp. 263--270. Springer, 2014.

\bibitem[Jarvis et~al.(2008)]{jarvis2008srtm}
Jarvis, A., Reuter, H.I., Nelson, A., and Guevara, E.
\newblock Hole-filled SRTM for the globe Version 4.
\newblock CGIAR-CSI, 2008.

\bibitem[Jetz et~al.(2019)]{jetz2019essential}
Jetz, W., McGeoch, M.A., Guralnick, R., et~al.
\newblock Essential biodiversity variables for mapping and monitoring species
  populations.
\newblock \emph{Nature Ecology \& Evolution}, 3:539--551, 2019.

\bibitem[Kipf \& Welling(2017)]{kipf2017semi}
Kipf, T.N. and Welling, M.
\newblock Semi-supervised classification with graph convolutional networks.
\newblock \emph{ICLR}, 2017.

\bibitem[Landau et~al.(2021)]{landau2021omnidirectional}
Landau, V.A., Shah, V.B., Anantharaman, R., and Hall, K.R.
\newblock Omniscape.jl: Software to compute omnidirectional landscape
  connectivity.
\newblock \emph{JOSS}, 6(57):2829, 2021.

\bibitem[Lek \& Gu{\'e}gan(1999)]{lek1996application}
Lek, S. and Gu{\'e}gan, J.F.
\newblock Artificial neural networks as a tool in ecological modelling.
\newblock \emph{Ecological Modelling}, 120(2-3):65--73, 1999.

\bibitem[Li et~al.(2018)]{li2018diffusion}
Li, Y., Yu, R., Shahabi, C., and Liu, Y.
\newblock Diffusion convolutional recurrent neural network: Data-driven traffic
  forecasting.
\newblock \emph{ICLR}, 2018.

\bibitem[McRae(2006)]{mcrae2006isolation}
McRae, B.H.
\newblock Isolation by resistance.
\newblock \emph{Evolution}, 60(8):1551--1561, 2006.

\bibitem[McRae et~al.(2008)]{mcrae2008using}
McRae, B.H., Dickson, B.G., Keitt, T.H., and Shah, V.B.
\newblock Using circuit theory to model connectivity in ecology, evolution, and
  conservation.
\newblock \emph{Ecology}, 89(10):2712--2724, 2008.

\bibitem[Nemhauser et~al.(1978)]{nemhauser1978analysis}
Nemhauser, G.L., Wolsey, L.A., and Fisher, M.L.
\newblock An analysis of approximations for maximizing submodular set
  functions---I.
\newblock \emph{Mathematical Programming}, 14(1):265--294, 1978.

\bibitem[Newmark(2002)]{newmark2002conserving}
Newmark, W.D.
\newblock Conserving biodiversity in East African forests.
\newblock \emph{Ecological Studies}, 155, Springer, 2002.

\bibitem[Pecl et~al.(2017)]{pecl2017biodiversity}
Pecl, G.T., Ara{\'u}jo, M.B., Bell, J.D., et~al.
\newblock Biodiversity redistribution under climate change: Impacts on
  ecosystems and human well-being.
\newblock \emph{Science}, 355(6332):eaai9214, 2017.

\bibitem[Phillips et~al.(2006)]{phillips2006maximum}
Phillips, S.J., Anderson, R.P., and Schapire, R.E.
\newblock Maximum entropy modeling of species geographic distributions.
\newblock \emph{Ecological Modelling}, 190(3-4):231--259, 2006.

\bibitem[Roberts et~al.(2017)]{roberts2017cross}
Roberts, D.R., Bahn, V., Ciuti, S., et~al.
\newblock Cross-validation strategies for data with temporal, spatial,
  hierarchical, or phylogenetic structure.
\newblock \emph{Ecography}, 40(8):913--929, 2017.

\bibitem[Strydom et~al.(2021)]{strydom2021graph}
Strydom, T., Catchen, M.D., Banville, F., et~al.
\newblock A roadmap towards predicting species interaction networks.
\newblock \emph{Nature Ecology \& Evolution}, 5:790--798, 2021.

\bibitem[Sullivan et~al.(2009)]{sullivan2009ebird}
Sullivan, B.L., Wood, C.L., Iliff, M.J., et~al.
\newblock eBird: A citizen-based bird observation network.
\newblock \emph{Biological Conservation}, 142(10):2282--2292, 2009.

\bibitem[Thuiller et~al.(2009)]{thuiller2009biomod}
Thuiller, W., Lafourcade, B., Engler, R., and Ara{\'u}jo, M.B.
\newblock BIOMOD---a platform for ensemble forecasting of species distributions.
\newblock \emph{Ecography}, 32(3):369--373, 2009.

\bibitem[Urban(2015)]{urban2015accelerating}
Urban, M.C.
\newblock Accelerating extinction risk from climate change.
\newblock \emph{Science}, 348(6234):571--573, 2015.

\bibitem[Veli{\v{c}}kovi{\'c} et~al.(2018)]{velickovic2018graph}
Veli{\v{c}}kovi{\'c}, P., Cucurull, G., Casanova, A., et~al.
\newblock Graph attention networks.
\newblock \emph{ICLR}, 2018.

\bibitem[Zanaga et~al.(2022)]{zanaga2022esa}
Zanaga, D., Van De~Kerchove, R., Daems, D., et~al.
\newblock ESA WorldCover 10m 2021 v200.
\newblock Zenodo, 2022.

\bibitem[Zurell et~al.(2016)]{zurell2016benchmarking}
Zurell, D., Thuiller, W., Pagel, J., et~al.
\newblock Benchmarking novel approaches for modelling species range dynamics.
\newblock \emph{Global Change Biology}, 22(8):2651--2664, 2016.

\end{thebibliography}

% ══════════════════════════════════════════════════════════════════════════════
\appendix
\section{Supplementary: Feature Importance Analysis}
\label{app:features}
% ══════════════════════════════════════════════════════════════════════════════

We analyze feature importance using integrated gradients
\citep{chen2020deep} applied to the feature encoder. Across all 847
species, the most influential variables are:

\begin{table}[H]
\centering
\caption{Top 10 features by mean integrated gradient magnitude across all
species.}
\label{tab:features}
\begin{tabular}{clc}
\toprule
\textbf{Rank} & \textbf{Feature} & \textbf{Importance} \\
\midrule
1  & Annual mean temperature (BIO1)      & 0.187 \\
2  & Annual precipitation (BIO12)        & 0.156 \\
3  & Elevation                           & 0.134 \\
4  & Tree cover fraction                 & 0.121 \\
5  & Temperature seasonality (BIO4)      & 0.098 \\
6  & NDVI mean                           & 0.089 \\
7  & Precipitation seasonality (BIO15)   & 0.076 \\
8  & Distance to water                   & 0.062 \\
9  & Slope                               & 0.048 \\
10 & Temperature annual range (BIO7)     & 0.041 \\
\bottomrule
\end{tabular}
\end{table}

Feature importance varies substantially across taxa: for amphibians,
precipitation variables dominate; for birds, temperature and NDVI are
most influential; for mammals, tree cover and elevation are primary
drivers.

\section{Supplementary: Hyperparameter Sensitivity}
\label{app:hyperparameters}

\begin{table}[H]
\centering
\caption{Sensitivity of HabitatNet AUC to key hyperparameters.}
\label{tab:hyperparams}
\begin{tabular}{lcccc}
\toprule
\textbf{Parameter} & \textbf{Values tested} & \textbf{Best} & \textbf{AUC range} \\
\midrule
Message-passing layers $L$ & 2, 4, 6, 8, 10 & 6 & 0.931--0.943 \\
Hidden dimension & 64, 128, 256 & 128 & 0.937--0.943 \\
Learning rate & $10^{-3}$, $3\times10^{-4}$, $10^{-4}$ & $3\times10^{-4}$ & 0.935--0.943 \\
Spatial loss $\lambda_1$ & 0.01, 0.05, 0.1, 0.5 & 0.1 & 0.938--0.943 \\
Connectivity loss $\lambda_2$ & 0.01, 0.05, 0.1 & 0.05 & 0.940--0.943 \\
\bottomrule
\end{tabular}
\end{table}

The model is most sensitive to the number of message-passing layers and
least sensitive to the connectivity loss weight, suggesting that the
spatial structure provides the dominant signal.

\end{document}
